\emph{The tests.} Let $J$ be the unit of the boundary-mode arrays, $J^{2} = -1$, central among the modes (on the continuum light ray the unit is $I_{4}$, and everything below holds with $I_{4}$ for $J$). For $c \in \Z$, $e^{2\pi Jc/24} = \cos\tfrac{2\pi c}{24} + J\sin\tfrac{2\pi c}{24} = 1$ exactly when $c/24 \in \Z$; likewise $e^{2\pi Jc/12} = (e^{2\pi Jc/24})^{2} = 1$ exactly when $12 \mid c$, and $e^{2\pi Jc/24} = \pm1$ exactly when $12 \mid c$. An operator implementing the grading is fixed only up to sign; normalising it by $P = +1$ on the sea, the one-window test says that $e^{2\pi Jc/24}P$ is $P$, the two-window test that it is $\pm P$, equivalently that its square, the evolution over two windows, is $1$ ($P^{2} = 1$, $P$ commuting with the central unit). The form $e^{2\pi Jc/24}P$ is the reading under which the tests are posed: each complex array carries $-\kappa/24$ over a window $2\pi/\kappa$ in the null ordering (Theorem \claimref{synthesis.null-ordering-of-counted-units-is-zeta-symmetric}(i), (ii)), and a real array at integer weights carries, in the zeta-regularised symmetric ordering, $\tfrac12\sum_{n \ge 1}\kappa n = \tfrac{\kappa}{2}\zeta(-1) = -\kappa/24$ (the zero mode carrying nothing), with excitation energies in $\kappa\Z$, so its one-window phase on excitations is $1$. \emph{Computed:} the finite part at $\varepsilon = 0$ of $\sum_{n \ge 0}(n + a)e^{-\varepsilon(n + a)}$ is $-\tfrac12B_{2}(a)$, by exact series inversion, so $\zeta(-1) = -\tfrac12B_{2}(1) = -\tfrac1{12}$ and a real array carries $-\kappa/24$; a real array at the half offset would carry $\tfrac{\kappa}{2}\zeta(-1, \tfrac12) = +\kappa/48$, not the null-ordered window value. That the boson's zero point equals the null-ordered window value of the continuum, as Theorem \claimref{synthesis.null-ordering-of-counted-units-is-zeta-symmetric} shows for the CAR arrays, is part of the additive reading, not derived. None of this is used in (i) to (vi), which are statements about $c$.

\emph{The two readings.} Premise \claimref{hawking.joint-j3-s-wave-reduction} (J3), in the past cone of the null-ordering result, identifies each $(\omega, \ell, m, s)$ channel of every four-dimensional massless field with the current boson of the s-wave theory, which is a current of the arrays; read so, the integer-weight channels add no array and only the CAR count enters $c$. The additive reading instead counts their null data as further real arrays. Premise (N$_{+}$) and the window reading are stated for CAR arrays, so neither places vector, scalar or graviton data in the sea; the additive reading is an extension of both. Clauses (i) to (v) are statements under the additive reading and (vi) under J3's.

(i) \emph{Fermions and vectors} are the cited clauses: Theorem \claimref{synthesis.slot-module-light-ray-arrays-are-weyl-components}(i) shows that every Lorentz-covariant mass term, Dirac or Majorana, acts from the right and enters only the equation fixing $\Lambda(k)\psi$ from the free data $\Lambda(m)\psi$, one complex channel per Weyl component; Theorem \claimref{synthesis.null-plane-photon-data-and-the-pair-currents-of-the-generators}(ii) shows that the flux of the Maxwell symplectic current through $N_{k}$ in the gauge $k\cdot A = 0$ contains the two transverse components alone, so an abelian gauge vector gives two real fields on each generator. \emph{Argued, not proved:} for the factors of a non-abelian algebra the count is that of the free field, component by component in the adjoint, the interaction entering, as the mass terms do for the spinor, only through the equations off the plane (light-front gauge); no stored clause covers it. Both carry integer or half-integer boost weight as recorded in Theorem \claimref{synthesis.light-ray-good-part-of-a-spinor-is-two-chirality-arrays}(ii), (vi): spinor data weight $\tfrac12$ (the complex arrays), vector and bivector data integer weight (the real arrays); clause (vi) of that theorem does not cover scalars.

\emph{Scalars.} Write $\nabla = \gamma^{\mu}\partial_{\mu}$ with the reciprocal frame $\gamma^{\mu}\cdot\gamma_{\nu} = \delta^{\mu}_{\nu}$, so $a\cdot\nabla = a^{\mu}\partial_{\mu}$ and $\nabla^{2} = \partial_{0}^{2} - \partial_{1}^{2} - \partial_{2}^{2} - \partial_{3}^{2}$. Let $\phi$ solve $\nabla^{2}\phi + V'(\phi) = 0$ and let $\phi_{1}, \phi_{2}$ solve the linearised equation $\nabla^{2}\phi_{i} + V''(\phi)\phi_{i} = 0$. For $\omega = \phi_{1}\nabla\phi_{2} - \phi_{2}\nabla\phi_{1}$, the Leibniz rule gives $\nabla\cdot\omega = (\nabla\phi_{1})\cdot(\nabla\phi_{2}) + \phi_{1}\nabla^{2}\phi_{2} - (\nabla\phi_{2})\cdot(\nabla\phi_{1}) - \phi_{2}\nabla^{2}\phi_{1} = \phi_{1}\nabla^{2}\phi_{2} - \phi_{2}\nabla^{2}\phi_{1}$ identically, and on the linearised solutions this is $-\phi_{1}V''\phi_{2} + \phi_{2}V''\phi_{1} = 0$: $\omega$ is conserved for every potential. The plane $N_{k} = \{x : \ip{x}{k} = 0\}$ has normal $k$, which is null and so tangent to it; the flux of $\omega$ through $N_{k}$ is $k\cdot\omega = \phi_{1}(k\cdot\nabla)\phi_{2} - \phi_{2}(k\cdot\nabla)\phi_{1}$, with $k\cdot\nabla = \partial_{0} + \partial_{3}$ the derivative along the generators. It contains no potential and no transverse derivative, and as for the vector field it is, per generator, the symplectic form $\int(\phi_{1}\partial_{u}\phi_{2} - \phi_{2}\partial_{u}\phi_{1})\,du$ of one real chiral boson in the affine parameter $u$. With $m\cdot\nabla = \partial_{0} - \partial_{3}$, $(k\cdot\nabla)(m\cdot\nabla) = \partial_{0}^{2} - \partial_{3}^{2}$, so $\nabla^{2} = (k\cdot\nabla)(m\cdot\nabla) - \partial_{1}^{2} - \partial_{2}^{2}$, and the field equation reads $(k\cdot\nabla)\bigl((m\cdot\nabla)\phi\bigr) = (\partial_{1}^{2} + \partial_{2}^{2})\phi - V'(\phi)$: on $N_{k}$ it is an ordinary differential equation along each generator for $(m\cdot\nabla)\phi$, with right side known from $\phi$ on $N_{k}$, so it fixes $(m\cdot\nabla)\phi$ there up to one boundary value per generator and constrains nothing else. The free data are $\phi$ on $N_{k}$, one real field per generator. \emph{Argued:} their weight under the boost of the generator is an integer (a scalar's values are unchanged by the boost, weight $0$, and its derivative along the generator has weight $1$), so a real scalar component is one real array; the stored weight clause covers vectors and bivectors only. \emph{Computed:} $Q(k) = Q(m) = 0$, $\ip{k}{m} = 2$, the divergence identity, the flux formula and the factorisation of $\nabla^{2}$ on polynomial fields in $\Cl(1,3)$; and, as a control, that the flux through a surface of constant time differs. Adding the counts gives $c = 16n_{g} + 2n_{V} + n_{s} + 2\epsilon$, where $\epsilon$ records whether a graviton's two helicities are counted as two further real arrays (J3's species, every four-dimensional massless field, include the graviton, so a graviton counted additively has $\epsilon = 1$; $\epsilon$ is left as a parameter).

(ii) Under the additive reading $2n_{V} = 24$ for $n_{V} = 12$, the dimension of $\su(3) \oplus \su(2) \oplus \uu(1)$ (Theorem \claimref{standard-model.connection-forced}). For dimension $d$, $24 \mid 2d$ if and only if $12 \mid d$ and $12 \mid 2d$ if and only if $6 \mid d$. \emph{Computed:} $d = 0, \dots, 39$; $\mathfrak{so}(10)$, $d = 45$, gives $90$ and fails both.

(iii) With $n_{V} = 12$, $n_{s} = 4n_{h}$: $c = 16n_{g} + 24 + 4n_{h} + 2\epsilon$. If $\epsilon = 1$, $16n_{g} + 24 + 4n_{h}$ is divisible by $4$, so $c \equiv 2 \pmod 4$ and $c$ is divisible neither by $12$ nor by $24$. If $\epsilon = 0$: $24 \mid 16n_{g} + 4n_{h}$ if and only if $6 \mid 4n_{g} + n_{h}$, that is $n_{h} \equiv -4n_{g} \equiv 2n_{g} \pmod 6$; and $12 \mid 16n_{g} + 4n_{h}$ if and only if $3 \mid 4n_{g} + n_{h}$, that is $n_{g} + n_{h} \equiv 0 \pmod 3$. The cases: $n_{h} = 0$ gives $3 \mid n_{g}$ in both; $n_{h} = 1$ is odd while $2n_{g}$ is even, so the one-window test never holds, and $n_{g} \equiv 2 \pmod 3$ for the other; $n_{h} = 2$ gives $2 \equiv 2n_{g} \pmod 6$, i.e. $n_{g} \equiv 1 \pmod 3$, and $n_{g} + 2 \equiv 0 \pmod 3$, the same. At $n_{g} = 2$, $n_{h} = 1$: $c = 60$, $c/24 = \tfrac52$ and $e^{2\pi J\cdot 5/2} = J^{10} = -1$ (the window evolution is $-P$); at $n_{g} = 3$, $n_{h} = 2$: $c = 80$, $80/12$ is not an integer; at $n_{g} = 3$, $n_{h} = 1$, $\epsilon = 1$: $c = 48 + 24 + 4 + 2 = 78$. The linear rung field has eight real components (Theorem \claimref{synthesis.linear-rung-field-has-a-neutral-normal-pair-in-the-scalon-sum}, its four $I_{4}$-complex modes per momentum), so $n_{s} = 8$, $n_{h} = 2$. \emph{Computed:} both equivalences on a full period ($n_{g} \bmod 3$, $n_{h} \bmod 6$, $\epsilon$) with two lifts, and the one-window rule on $0 \le n_{g}, n_{h} < 40$; the pass lists for $1 \le n_{g} \le 9$ at $n_{h} = 0, 1, 2$; the values $J^{10} = -1$, $J^{20} = 1$ with $J = \gamma_{1}\gamma_{2}$ of $\Cl(1,3)$ standing for the unit.

(iv) Let $g = \gcd(24, f)$. Then $24 \mid fn_{g}$ if and only if $(24/g) \mid (f/g)n_{g}$, and since $24/g$ and $f/g$ are coprime, if and only if $(24/g) \mid n_{g}$; the same with $12$. $24/g = 3$ exactly when $g = 8$, that is $8 \mid f$ and $3 \nmid f$, that is $f \equiv 8$ or $16 \pmod{24}$. For $f = 15$, $g = 3$ and $24/g = 8$, $12/\gcd(12,15) = 4$. That fifteen components are not a CPT halving of the slot module is Theorem \claimref{synthesis.slot-module-light-ray-arrays-are-weyl-components}(ii). \emph{Computed:} every $f \le 48$ against every $n_{g} \le 96$.

(v) If $b = 2f$ then $c = f + b = 3f = 48n_{g}$ for $f = 16n_{g}$, a multiple of $24$. In the notation of Theorem \claimref{synthesis.attachment-modules-give-sakharov-weights-cancelling-at-three-separate-generations}(iv) the quartic coefficient is $\tfrac12(2N_{V} + N_{S} - 2N_{W})$ times a positive integral of the cut-off, with $2N_{V} + N_{S}$ the bosonic helicity states and $2N_{W}$ the fermionic ones; in the count of (i) with $\epsilon = 0$, $b = 2N_{V} + N_{S}$ and $f = N_{W}$, so its vanishing is $b = 2f$. At $N_{g} = 3$, $N_{V} = 12$, $N_{W} = 48$: $N_{S} = 96 - 24 = 72$ and $c = 48 + 24 + 72 = 144 = 6 \cdot 24$; one scalar fewer gives $143$, which fails.

(vi) Under J3's reading $b = 0$ and $c = f = 16n_{g}$; by (iv) with $f = 16$, $\gcd(24, 16) = 8$ and $\gcd(12, 16) = 4$, so both tests hold exactly when $3 \mid n_{g}$. \emph{Computed:} $c = 16n_{g}$ passes exactly $n_{g} = 3, 6, 9$ for $1 \le n_{g} \le 9$, under both tests.

\emph{What is proved and what is not.} Proved: (i) to (vi) as statements about $c$ for the listed contents under the stated reading, the scalar null-data clause in geometric algebra; argued: the non-abelian vector count and the integer weight of scalar data. Not proved and not claimed: that the recurrence of the light-ray sea is a premise of the series, or which fields form the sea. Premise (N$_{+}$) is stated for the CAR modes of the light rays, so the store puts no vector or scalar null data in the null ground state, and J3 reads the integer-weight channels as currents of the arrays; whether they should instead be counted additively, and whether the graviton's helicities then count, are premise decisions. Neither test separates $n_{g} = 3$ from $6$ or $9$. Argued: under spontaneous breaking the count is unchanged, three Goldstone components becoming the longitudinal data of three massive vectors ($2 \cdot 9 + 3 \cdot 3 + 1 = 28 = 2 \cdot 12 + 4$). Transport-absent.
